\subsection{Tensor Products of Semisimple Commutative Algebras}

THEOREM 3. --- Let \(\mathbf{Z}_1\) and \(\mathbf{Z}_2\) be semisimple commutative algebras over \(\mathbf{K}\) . The radical of the ring \(\mathbf{Z}_1 \otimes \mathbf{Z}_2\) is equal to the set of its nilpotent elements.

Let us first treat the case when \(Z_{1}\) and \(Z_{2}\) are extensions \(L_{1}\) and \(L_{2}\) of the field K. By interchanging \(L_{1}\) and \(L_{2}\) if necessary, we reduce to the case when the transcendent degree of \(L_{1}\) over K is less than that of \(L_{2}\) over K.

Choose an algebraic closure \(\Omega\) of \(\mathrm{L}_2\) ; by Corollary 1 of Theorem 5 of V, §14, No.~6, p.~114, we may assume that \(\mathrm{L}_1\) is a subextension of \(\Omega\) .

\begin{enumerate}
\def\labelenumi{\Alph{enumi})}
\item
  Let us first prove that the radical of \(L_{1} \otimes L_{2}\) is contained in that of \(L_{1} \otimes \Omega\) . Set \(\mathfrak{a} = \Re(L_{1} \otimes L_{2})(L_{1} \otimes \Omega)\) ; it is an ideal of the commutative ring \(L_{1} \otimes \Omega\) , and we must prove that a is contained in the radical of \(L_{1} \otimes \Omega\) . In other words (VIII, p.~156, Theorem 1), we must prove that for \(x \in a\) , the element \(1 + x\) is invertible in \(L_{1} \otimes \Omega\) . Now, since \(\Omega\) is an algebraic extension of \(L_{2}\) , there exists an extension \(L_{3}\) of \(L_{2}\) , of finite degree, such that x belongs to \(\Re(L_{1} \otimes L_{2})(L_{1} \otimes L_{3})\) . It obviously suffices to prove that \(1 + x\) is invertible in \(L_{1} \otimes L_{3}\) . Now, \(C = L_{1} \otimes L_{3}\) is a finitely generated module over the ring \(B = L_{1} \otimes L_{2}\) . By the corollary of VIII, p.~175, we have \(\Re(B)C \subset \Re(C)\) , so x belongs to the radical of C and \(1 + x\) is invertible in C.
\item
  Let us prove that the radical of \(L_{1} \otimes \Omega\) consists of nilpotent elements. Denote the characteristic exponent of K by p and the relative p-radical (that is, purely inseparable) closure of K in \(\Omega\) (V, §5, No.~2, p.~25) by P; the field P is perfect. Since P is an algebraic extension of K, we have \(\mathrm{L}_{1}(\mathrm{P}) = \mathrm{L}_{1}[\mathrm{P}]\) (V, §3, No.~2, p.~18, Corollary 1). Let b be the kernel of the canonical homomorphism from \(L_{1} \otimes P\) to the field \(P_{1} = L_{1}[P]\) . Let \(x \in b\) ; there exist elements \(y_{1}, \ldots, y_{n}\) of \(L_{1}\) and elements \(z_{1}, \ldots, z_{n}\) of P such that \(x = \sum_{i=1}^{n} y_{i} \otimes z_{i}\) and \(\sum_{i=1}^{n} y_{i} z_{i} = 0\) . Since P is p-radical over K, there exists a power q of p such that \(z_{1}^{q}, \ldots, z_{n}^{q}\) belong to K. We then have
\end{enumerate}

\[
x ^ {q} = \sum_ {i = 1} ^ {n} y _ {i} ^ {q} \otimes z _ {i} ^ {q} = \sum_ {i = 1} ^ {n} y _ {i} ^ {q} z _ {i} ^ {q} \otimes 1 = \left(\sum_ {i = 1} ^ {n} y _ {i} z _ {i}\right) ^ {q} \otimes 1 = 0.
\]

So b consists of nilpotent elements.

Set \(c = b \otimes_{P} \Omega\) ; it is the kernel of the canonical homomorphism from \((\mathrm{L}_{1} \otimes \mathrm{P}) \otimes_{\mathrm{P}} \Omega\) to \(P_{1} \otimes_{P} \Omega\) , and it consists of nilpotent elements by the above. Now, \(\Omega\) is an algebraically closed extension of P, and \(P_{1}\) is a subextension of \(\Omega\) . Since the field P is perfect, \(P_{1}\) is a separable extension of P (V, §15, No.~5, p.~125, Theorem 3). By Theorem 4 of V, p.~120, the intersection of the maximal ideals of the commutative ring \(P_{1} \otimes_{P} \Omega\) is reduced to 0. In other words, the ring \(P_{1} \otimes_{P} \Omega\) , which is isomorphic to \(((L_{1} \otimes P) \otimes_{P} \Omega)/c\) , is without radical. This proves (VIII, p.~155, Proposition 5) that c contains the radical of the ring \((\mathrm{L}_{1} \otimes \mathrm{P}) \otimes_{\mathrm{P}} \Omega\) . Now, this ring is isomorphic to \(L_{1} \otimes \Omega\) , and c consists of nilpotent elements. Therefore, the radical of \(L_{1} \otimes \Omega\) consists of nilpotent elements.

\begin{enumerate}
\def\labelenumi{\Alph{enumi})}
\setcounter{enumi}{2}
\tightlist
\item
  End of the proof in the specific case. By A) and B), the radical \(\mathfrak{r}\) of \(\mathrm{L}_1\otimes \mathrm{L}_2\) is contained in the set \(\mathfrak{n}\) of nilpotent elements of this commutative ring. We moreover know that \(\mathfrak{n}\) is contained in \(\mathfrak{r}\) (VIII, p.~157, Remark 2).
\end{enumerate}

Now consider the general case. Since a semisimple commutative K-algebra is the product of finitely many extensions of the field K (VIII, p.~137, Proposition 3) and the radical of a product of rings is the product of the radicals (VIII, p.~156, Corollary 3), the radical of \(Z_{1} \otimes Z_{2}\) is the set of nilpotent elements of this ring.

